\scanpage{7}
\[
A_n=A/(F^{\unsure{n+1}},V^{\unsure{n+1}})
\]
\[
I_n=A_n\quad\text{regardé comme \uline{bimodule} sur }A
\]
\[
\uline{I=\underset{n}{\varinjlim}\,I_n}
\]
\diagram{\begin{tikzcd}I_n\arrow[r,dash,"p"]&I_{n+1}\end{tikzcd}}
$D:M\rightsquigarrow\operatorname{Hom}(M,\unsure{I|I_n})$ : foncteur dualisant

\uline{Attention}, à priori un module à g. est transformé en module à droite, et vice versa. Mais on peut utiliser l'involution $s$ de $A$, échangeant $F$ et $V$, pour passer de l'un à l'autre.

\deleted{$D(M)$} $A_n$ est fini sur $W_n$

Un autre foncteur dualisant est obtenu en faisant (si $M$ tué par $F^n,V^n$)
\[
D':M\rightsquigarrow\operatorname{Hom}_{W_n\text{-modules}}(M,W_n)
\]
et définissant sur $D'(M)$ la structure de $A$-module à droite par ``\unsure{transposé}''
\inserted{$D'(A_n)\simeq\operatorname{Hom}(A_n,W_n)$ structure de bimodule}
\[
\operatorname{Hom}_{W_n}(M,W_n)=\operatorname{Hom}_{A_n}\bigl(M,\underbrace{\operatorname{Hom}_{W_n}(A_n,W_n)}_{\substack{\text{structure}\\\text{de bimodule}\\\text{naturelle}}}\bigr)
\]
Je dis qu'on a un iso
\[
\uline{I_n=A_n\simeq\operatorname{Hom}_{W_n}(A_n,W_n)}
\]
\deleted{$\unsure{uF=\pi uF}$}
\mspage{6}
